% 11_body.tex — Day3 산출물 1/5 (⑪) 변경요청서·영향 분석·CCB 의결서 (빈 템플릿 / 완성 예시 공용 본문)
% 드라이버: 11_변경요청서_빈템플릿.tex (\wbblanktrue) / 11_변경요청서_완성예시.tex (\wbblankfalse)
% 내용 출처: PM트랙/Day3_워크북.md §1.3(항목 가이드) §1.4(빈 템플릿) §1.5(온담 P1 완성 예시본)
\documentclass[10pt,a4paper]{article}
\usepackage[a4paper,margin=16mm,top=14mm,bottom=20mm]{geometry}
\def\DocNum{1}
\def\DocTitle{변경요청서 · 영향 분석 · CCB 의결서}
\def\DocSub{⑪ Change Request, Impact Analysis \& CCB Decision Record}
\def\DocVariant{\B{빈 템플릿}{온담 P1 완성 예시}}
\usepackage{wbstyle}
\begin{document}

\docheader
 {\Bg{[정식 명칭과 코드]}{차세대 주문·재고 통합 플랫폼 구축 (ONE ONDAM P1)}}
 {\Bg{[CR 번호 범위 / 데이터 일자 / CCB 일자]}{CR-08~11 / 2026-09-30 / CCB 제10차 2026-10-15}}
 {\Bg{[P1 PMO 요구사항 담당 — 접수자]}{P1 PMO 요구사항 담당 (접수), P1 PM (영향 분석)\rd{7mm}}}
 {\Bg{[CCB 위원장 실명]}{정미란 CFO (CCB 위원장)\rd{7mm}}}

% ------------------------------------------------------------------ 0 처리 경로 규칙
\secbar[45mm]{0. 처리 경로 규칙 (변경 관리 계획에서 인용)}
\guide{변경요청서를 채우기 전에 이 프로젝트의 처리 경로 규칙을 문서 앞에 적는다. 규칙은 두 개의 문을 갖는다 — \textbf{허용범위는 크기의 기준, 헌장 문언은 권한의 기준.} 허용범위 안이라도 헌장 §5·§12·§13의 문장을 바꾸면 CCB이고, 타 프로젝트·운영 조직에 영향이 있으면 운영위 보고를 병행한다. 승인된 변경은 건별로 기준선을 고치지 않고 BL 번호를 올려 일괄 갱신한다.}
\begin{fieldbox}
\ifwbblank
\g{① [모든 축 허용범위 내 · 헌장 문언 불변 · 타 프로젝트 영향 없음] → [PM 전결 + CCB 사후 보고]}\par
\g{② [헌장 문언 변경 또는 어느 축 초과] → [CCB; 수행사 FP 영향 시 변경협의체 합의 선행]}\par
\g{③ [관리예비비 필요] → [스폰서 결재]}\qquad \g{④ [타 프로젝트·운영 영향] → [운영위 보고 병행]}\par
\g{⑤ [릴리스 이연] → [헌장 §13 인용, CCB]}\qquad \g{⑥ [승인 후 n영업일 내 BL 번호 일괄 갱신]}
\else
\small \textbf{변경 관리 계획 v1.0 §3 (G1 확정 [추가 설정])} — ① 6축 모두 허용범위 내, 헌장 §5·§12·§13 문언 불변, 타 프로젝트·운영 영향 없음 → P1 PM 전결(RACI 16행), 다음 CCB 사후 보고. ② 헌장 문언 변경 또는 어느 축 초과 → CCB 승인(RACI 17행); 수행사 FP 영향이 있으면 변경협의체 합의 선행(RACI 18행). ③ 관리예비비 필요 → 스폰서 결재(RACI 20행). ④ 타 프로젝트·운영 조직 영향 → 프로그램 운영위 보고 병행. ⑤ 릴리스 이연(R5) → 헌장 §13 인용, CCB. ⑥ 승인된 변경은 CCB 후 5영업일 내 BL 번호를 올려 기준선 문서 일괄 갱신.
\fi
\end{fieldbox}

% ================================================================== 변경요청서 (항목 1~7) — 빈 템플릿 1부 / 완성 예시 4건
% 공용 매크로: \crsheet{번호·제목}{항목1 표 4셀}{항목2 내용}{항목3 표 3셀}{항목4~6 표 행들}{항목7 표 행들}{판정 문장}
\newcommand{\crsheet}[7]{%
\secbar[60mm]{변경요청서 — #1}
\par\vspace{1mm}\noindent\textbf{\small 1. 변경 식별}\quad\g{CR 번호(RTM 변경 이력 열과 같은 체계)·제목·요청자(실명·S-ID)·접수 경로·접수일·동결 회차 관계. PM은 요청자가 아니라 접수자다. 기각되어도 번호를 회수하지 않는다.}\par\vspace{0.5mm}
{\footnotesize
\begin{tabularx}{\linewidth}{|L{24mm}|Y|L{22mm}|L{40mm}|}\hline
\hd{항목} & \hd{내용} & \hd{항목} & \hd{내용} \\ \hline
#2
\end{tabularx}}
\par\vspace{2mm}\noindent\textbf{\small 2. 변경 내용과 사유}\quad\g{“현재 기준선 문장(문서·항목 인용) → 요청 문장” 두 줄 + 사유(판정 가능한 사실) + 안 바꾸면 + 대안 3개 이상. CCB가 “예/아니오”가 아니라 “어느 것”을 고르게 한다.}\par\vspace{0.5mm}
\begin{fieldbox}\small #3\end{fieldbox}
\par\vspace{2mm}\noindent\textbf{\small 3. 분류·긴급도}\quad\g{유형(범위 추가/축소/요구 변경/거버넌스/오류 정정/릴리스 이연)과 긴급도의 근거 일자. “긴급”은 규제·장애·계약 시한이 있을 때만.}\par\vspace{0.5mm}
{\footnotesize
\begin{tabularx}{\linewidth}{|L{44mm}|Y|L{50mm}|}\hline
\hd{유형} & \hd{긴급도 · 근거 일자} & \hd{관련 R- / IS-} \\ \hline
#4
\end{tabularx}}
\par\needspace{40mm}\vspace{2mm}\noindent\textbf{\small 4~6. 6축 영향 분석 — 범위·일정 / 원가·품질 / 리스크·계약/편익}\quad\g{축마다 숫자에 출처. 일정은 TF로 말한다(“A12 TF 10 → +5 → TF 5, 임계경로 불변”). 원가는 산식과 자금원·잔액. 품질은 TC 번호와 RTM 갱신 행 수. 범위 축소는 FP를 음수로. 타 프로젝트 영향은 운영위 보고 대상 표시.}\par\vspace{0.5mm}
{\footnotesize
\begin{tabularx}{\linewidth}{|L{14mm}|Y|L{38mm}|}\hline
\hd{축} & \hd{영향} & \hd{근거} \\ \hline
#5
\end{tabularx}}
\par\needspace{45mm}\vspace{2mm}\noindent\textbf{\small 7. 허용범위 판정 · 처리 경로}\quad\g{격자(Day2 ⑩ 항목 7)의 축별 “내/초과”와 숫자, 누적(FP 누적 변동·컨틴전시 소진)과의 관계. 허용범위 내라도 헌장 문언을 바꾸면 CCB.}\par\vspace{0.5mm}
{\footnotesize
\begin{tabularx}{\linewidth}{|L{16mm}|Y|L{40mm}|L{42mm}|}\hline
\hd{축} & \hd{값} & \hd{허용범위 (격자)} & \hd{판정} \\ \hline
#6
\multicolumn{4}{|>{\raggedright\arraybackslash}p{\dimexpr\linewidth-2\tabcolsep-2\arrayrulewidth\relax}|}{\rule{0pt}{4mm}\small #7} \\ \hline
\end{tabularx}}\par}

\ifwbblank
% ---------------- 빈 템플릿: 1부
\crsheet{CR-[nn] [제목]}%
{\rd{9mm}CR 번호 / 제목 & \g{[CR-nn / 한 줄 제목]} & 접수일 / 동결 관계 & \g{[일자 / FZ-nn 접수 마감 전·후]} \\ \hline
\rd{9mm}요청자 / 접수 경로 & \g{[실명·S-ID / 협의회 서면·검토회·헌장 수정 요청 n번·감리·내부 발견]} & 접수자 & \g{[P1 PM / PMO]} \\ \hline}%
{\g{현재 기준선 문장 (문서·항목 인용):}\lines{2}{7mm}\g{요청 문장:}\lines{2}{7mm}\g{사유 (판정 가능한 사실 포함):}\lines{3}{7mm}\g{안 바꾸면 (편익·리스크·계약):}\lines{2}{7mm}\g{대안: (a) [요청대로] (b) [축소] (c) [이연/파라미터화] (d) [기각]}\lines{3}{7mm}}%
{\rd{11mm}\g{[범위 추가 / 축소 / 요구 변경 / 거버넌스 / 오류 정정 / 릴리스 이연]} & \g{[즉시 / 다음 동결 FZ-nn / 다음 CCB / 게이트 — 근거 일자]} & \g{[R-nn, IS-nn]} \\ \hline}%
{\rd{16mm}범위 & \g{WBS [코드] / RTM StRS [ ] SyRS [ ] SRS [ ] / TC [ ] / FP Δ [±n] (근거) / 누적 FP 변동 [ ]} & \g{[FP 산정서·WBS 사전 카드]} \\ \hline
\rd{16mm}일정 & \g{활동 [A-nn] 스프린트 [ ] / 기간 Δ [ ]일 / TF [전]→[후] / 임계경로 [불변/변동] / 마일스톤 [ ]} & \g{[CPM 계산표 행]} \\ \hline
\rd{16mm}원가 & \g{Δ [ ]억 = [산식] / 자금원 [기준선 내 / 컨틴전시 배정·미배정 풀 / 관리예비비 / 감액] / 잔액 [ ] / EAC 영향 [ ]} & \g{[원가 기준선 항목]} \\ \hline
\rd{16mm}품질 & \g{TC [+/−/재작성 n] / RTM 갱신 [행 수] / 결함밀도 분모 [FP] / 심각도 1·2 조건 [ ]} & \g{[RTM 시험 열]} \\ \hline
\rd{16mm}리스크 & \g{[R-nn P·영향·잔여 EMV 전→후] / 신규 [ ] / 잔여 합 [ ] vs 8억} & \g{[⑬ 항목 4]} \\ \hline
\rd{16mm}계약/편익 & \g{변경 대가 [증액/감액 억] / 계약 서면 [필요·대상] / 편익 [B-nn 영향] / 타 프로젝트 [ ]} & \g{[변경협의체 합의서]} \\ \hline}%
{\rd{7mm}범위 & \g{[건수 / FP Δ, 누적]} & \g{±35건 / FP ±372} & \\ \hline
\rd{7mm}일정 & \g{[TF 전→후, 마일스톤]} & \g{+15영업일 / 창 0} & \\ \hline
\rd{7mm}원가 & \g{[Δ, 컨틴전시 누적 소진]} & \g{10.2} & \\ \hline
\rd{7mm}품질 & \g{[TC, 커버리지]} & \g{0.4건/FP · 미해결 0} & \\ \hline
\rd{7mm}리스크 & \g{[잔여 EMV 합]} & \g{8억 / 단일 2억} & \\ \hline}%
{\g{헌장 문언 변경: [예/아니오] \quad 타 프로젝트·운영 영향: [예/아니오] \quad 처리 경로: [PM 전결+CCB 사후 / 변경협의체→CCB / 스폰서 / 운영위 보고 병행] \quad 계약 서면: [ ]}\lines{2}{7mm}}
\else
% ---------------- 완성 예시: CR-08
\crsheet{CR-08 파일럿 12개점 컷오버 전 임시 재고 동기화 연동 (범위 추가)}%
{CR 번호 / 제목 & CR-08 / 파일럿 12개점(P3 지정) 컷오버 전 OD-POS↔신 POS 임시 재고 동기화 연동 및 파일럿 배포 도구 & 접수일 / 동결 관계 & 2026-09-16 (FZ-10 접수 마감 10-14 전) \\ \hline
요청자 / 접수 경로 & 정미란 스폰서 (헌장 v1.1 수정 요청 1번 — Day2 §6, G1에서 “FP 산정 후 CCB 상정” 조건부 유보) / 발견자 P1 PM, FP 산정서 부록 v1.1 (이재현, 2026-09-10) & 접수자 & P1 PMO 요구사항 담당 \\ \hline}%
{\textbf{현재 기준선 문장:} 범위 기술서 v1.0 §2 산출물 ① “파일럿 12개점(P3 지정) 컷오버 전 배포” / WBS 1.1.7 “파일럿 12개점 배포·임시 연동 — 공동(P3 PM, 수행 POS팀), R4 S27–S28” / 헌장 v1.1 §5 포함 범위에 임시 연동 \textbf{없음} / FP 기준선 12,400에 미포함.\par
\textbf{요청 문장:} 헌장 §5 산출물 ①에 “파일럿 12개점은 컷오버 전 2027-01-04~02-25 기간 OD-POS와 신 POS를 병행 운영하며, 그 기간의 매장 재고를 양 시스템 간 동기화하는 임시 연동 2본(신 POS → OD-POS 판매 차감, OD-POS → 신 POS 야간 재고 정정)과 파일럿 배포·회수 도구를 P1이 인도한다”를 추가. WBS 1.1.7 산출물에 위 2본·도구 추가, FP +90.\par
\textbf{사유:} P3가 2026-08-28 파일럿 12개점(직영 9·가맹 3, 구형 기종 2점 포함)을 확정했다. 파일럿 기간 중 매장은 신 POS로 판매하지만 본사 재고·발주·SAP 정산은 OD-POS 원장으로 돌아간다(컷오버 전). 동기화가 없으면 파일럿 매장의 재고가 두 원장에서 갈라져 발주 오류가 나고, 2019년 스마트오더가 “매장 운영 부하”로 중단된 경로를 12개점에서 재현한다. FP 산정서 부록 v1.1: 연동 2본 68 FP + 배포·회수 도구 22 FP = 90 FP.\par
\textbf{안 바꾸면:} 파일럿 무효(G3 통과 조건 “파일럿 2주 운영 중 심각도 1·2 결함 0건” 판정 불가) 또는 2027-01 변경 대가 협상(수행사 우위, 코드 프리즈 중 개발).\par
\textbf{대안:} (a) 요청대로 R4 S27–S28(원안 WBS)에서 개발 — 코드 프리즈(2027-01-04~) 위반 \quad (b) \textbf{R4 잔여 A12(S25–S26, 2026-12-07~2027-01-01)로 앞당겨 개발, 파일럿 배포 S27} \quad (c) 파일럿 매장을 직영 9점으로 줄이고 수기 대사 — P3 결정 사항, 가맹 3점 제외 시 협상 지렛대 상실(R-15) \quad (d) 기각 — 파일럿 취소, G3 조건 개정 필요.}%
{범위 추가 (헌장 §5 문언 변경) & 다음 CCB 2026-10-15 — A12 착수 2026-12-07 전 FP 확정·수행사 자원 배정 필요, FZ-11(11-25) 동결에 StRS-095 포함 & R-13(구형 기종 파일럿), R-15(선행 릴리스 결함), R-19(신규: R4 잔여 압축), IS-04(슈퍼유저 이탈 — 파일럿 매장 재선정) \\ \hline}%
{범위 & WBS 1.1.7 산출물 +3(연동 2본·도구). RTM: StRS-095 신규(HR-03·HR-05 하위, “파일럿 병행 기간 재고 정합 유지”) → SyRS 3 / SRS 6 / TC-POS-181~192(12건). \textbf{FP +90} (68 + 22). 누적 FP 변동: +22(CR-03) +58(CR-06) +40(CR-07) +90 = \textbf{+210} (±372 이내, 56\%) & FP 산정서 부록 v1.1 §2; WBS 사전 1.1.7 카드(G1 작성분) 추정 “FP 0 — 배포 기록만” → 카드 갱신 \\ \hline
일정 & 원안 1.1.7 R4 S27–S28은 코드 프리즈 구간 → \textbf{A12(구현 R4 잔여, S25–S26, TF 10)에 편입}. A12 기간 20 → 25영업일(+5, FP 90 ≈ 스프린트 0.15회분이나 OD-POS 측 개발자 1명 병목으로 5일), \textbf{TF 10 → 5}, 임계경로 불변. 파일럿 배포는 S27(2027-01-04~17) 유지. A16(UAT) 시작 불변 & Day2 §3.5 CPM 표 A12(ES 240, TF 10, FF 10); R-06 소유자 박세연 확인 2026-09-14 \\ \hline
원가 & \textbf{+0.495억} = 90 × 0.0055. 자금원: 컨틴전시 \textbf{미배정 풀}(5.60 중 CR-03 0.12·CR-06 0.32·CR-07 0.22 기승인 0.66 차감 후 4.94 → 4.445). 발주사 공수 +0.3인월(파일럿 운영 P3·매장운영, P1 예산 외). EAC4 상향식 ETC에 포함(⑫ 항목 5) & Day2 §4.5 항목 6 미배정 풀 정의(“신규 식별 리스크·재작업”) — 범위 누락의 사후 발견은 이 용도에 해당 \\ \hline
품질 & TC +12(TC-POS-181~192), 정합 검증(양 원장 일치율 100\%) 케이스 포함. RTM 커버리지: StRS 214 → 215, 하위 분해·시험 연결 100\% 유지 조건. 결함밀도 분모 12,400 → 12,490(BL-02 시). 임시 연동은 컷오버 후 \textbf{폐기}(회수 절차 TC 2건) — 폐기 안 하면 운영 문서 47본에 유령 인터페이스 & RTM v1.1 §10 집계; 감리 G2 기준 \\ \hline
리스크 & \textbf{R-19}(⑬ v1.2에 등재, ⑫ 성과 분석에서 식별): R4 잔여(A12) 압축 → 옴니 Must·B2B EDI 시험 축소 → 통합시험 결함 유출, P 40 × 1.5 = 0.60(컨틴전시) — CR-08이 그 원인의 하나이며 승인 시 P 40 유지 확인. R-13 영향 불변(구형 기종 2점 파일럿 포함은 원안 유지). R-06 잔여 P 10 유지(OD-POS 측 개발자 투입 2주 확보). 잔여 위협 EMV 합 \textbf{7.26}(⑬, R-19 포함) 유지 — 8억 이내, 여유 0.74 & ⑬ 항목 4 \\ \hline
계약/편익 & 변경 대가 \textbf{증액 0.495억}, 변경협의체 2026-09-30 합의(FP 90 확인, 대기 없음). 계약 변경 서면(변경 합의서 6호)은 CR-09·CR-11과 묶어 10-23. 편익 영향 없음(B-ID 불변). 타 프로젝트: P3(파일럿 운영·교육 일정) 협의 완료 09-22 & 변경협의 합의서 2026-09-30 \\ \hline}%
{범위 & 건수 +1 / FP +90, 누적 +210 & ±35건 / FP ±372 & 내 \\ \hline
일정 & A12 TF 10 → 5, G4·컷오버 창 불변 & +15영업일, 창 0 & 내 \\ \hline
원가 & +0.495, 컨틴전시 누적 소진 1.41 → 1.905 & 10.2 & 내 \\ \hline
품질 & TC +12, 커버리지 100\% 유지 & 0.4건/FP · 미해결 0 & 내 \\ \hline
리스크 & 잔여 EMV 합 7.26 (R-19 포함) & 8억 / 단일 2억 & 내 (여유 0.74 — 경고) \\ \hline}%
{헌장 §5 문언 변경: \textbf{예} → \textbf{CCB 승인 필요}(허용범위 내이나 규칙 ②). 타 프로젝트 영향: P3 협의 완료, 운영위 보고 불요. 계약 서면: 변경 합의서 6호(10-23).}

% ---------------- 완성 예시: CR-09
\clearpage
\crsheet{CR-09 가맹 발주 마감 20:00 → 22:00 (요구 변경)}%
{CR 번호 / 제목 & CR-09 / 가맹 발주 포털 발주 마감 시각 20:00 → 22:00 변경 & 접수일 / 동결 관계 & 2026-09-17 협의회 서면 (FZ-10 접수 마감 10-14 전, 동결 10-28) \\ \hline
요청자 / 접수 경로 & S-10 서정필 가맹점주협의회장 / 협의회 공문 제2026-14호(오정근 경유) & 접수자 & P1 PMO \\ \hline}%
{\textbf{현재 기준선 문장:} RTM StRS-073 “발주 접수부터 출고 확정까지 12시간 이내가 되도록 마감·출고 확정 흐름이 설계된다” → FRN-FLOW-01 §2 “발주 마감 20:00, 이천센터 피킹 21:00 착수, 출고 확정 익일 08:00(12h)”. 범위 기술서 §2 산출물 ④ “발주 리드타임 12시간 설계”.\par
\textbf{요청 문장:} “발주 마감 22:00”. 협의회 사유서: 저녁 피크(19~21시) 마감 후 재고를 확인하고 발주하는 매장이 388점 중 약 210점(54\%, 협의회 자체 조사 [추가 설정])이며, 20:00 마감이면 이들은 전날 재고 기준으로 발주해 과잉·결품이 반복된다. 현행 가맹 발주 웹의 마감이 23:00이므로 20:00은 “후퇴”로 받아들여진다.\par
\textbf{안 바꾸면:} 협의회의 포털 사용성 검수(범위 기술서 §4 ④-③, 점주 5명) 거부 가능성, 부속서 잔여 112점 동의 협상(오정근)에 악재(R-02 P 상승).\par
\textbf{대안:} (a) 22:00 마감·피킹 23:00·출고 확정 익일 10:00 — 12h 유지, 센터 야간조 2시간 연장(물류본부, 대성 3PL 배송 출발 시각 변경) \quad (b) 22:00 마감·출고 확정 익일 12:00 — 14h, B-08 목표 재정의 \quad (c) \textbf{마감 시각을 운영 파라미터(매장군별 설정값)로 구현하고, 초기 운용은 20:00, 22:00 전환은 물류 야간조 편성 확정 후 CCB 재상정} \quad (d) 기각 — 20:00 유지, 협의회에 사유 회신.}%
{요구 변경 (StRS-073 하위 SRS 4건) & 다음 동결 FZ-10(10-28) — R3 포털 선행 릴리스(S23–S24, 2026-11-23~12-04)에 포함하려면 10-28 전 확정; 협의회 회신 기한 10-31(협의회 정기총회) & R-02(부속서 동의), R-16(폴백 배치 21:00·03:00과 충돌), IS-09 \\ \hline}%
{범위 & WBS 1.4.2(마감·출고 확정 흐름), 1.5.4(가맹발주↔SAP 출고 확정 인터페이스 IF-SAP-12 시각 파라미터). RTM StRS-073 변경 이력, SRS-413·414·416·417 수정(4건), SRS 신규 0. \textbf{FP +35}(파라미터화 12 + 마감 배치 재구성 15 + SAP 인터페이스 시각 변경 8). 누적 +245 & FP 산정서 부록 v1.1 §3; WBS 사전 1.4.2 카드 \\ \hline
일정 & 대안 (c) 파라미터 구현: A10(구현 R3, \textbf{TF 0, 현재 −10} — ⑫ 참조) 안의 S22(10-26~11-08)에서 15 FP, 마감 배치·SAP 시각은 R3 잔여 S23–S24 → \textbf{임계경로 A10에 +1영업일 부담} → 회복 계획(⑫ 대안)에 포함. 대안 (a)(b)는 물류 야간조 편성 전제 → 운용 전환은 R5(2027-03) 이후 & Day2 §3.5 A10 ES 180 / 실제 착수 작업일 190 \\ \hline
원가 & \textbf{+0.19억} = 35 × 0.0055. 자금원: 컨틴전시 미배정 풀. 대안 (a)는 P1 예산 외 운영비(야간조 인건비 연 약 0.9억 [추가 설정], 물류본부) & — \\ \hline
품질 & TC-FRN-040~048 중 3건 재작성, TC +8(TC-FRN-049~056: 매장군별 파라미터·전환 시험). UAT-FRN-039 수정. RTM 커버리지 불변 & RTM v1.1 시험 열 \\ \hline
리스크 & \textbf{R-16} 폴백 배치 부하(21:00 야간 배치 1회차와 22:00 마감분 피킹 충돌) P 30 → 45, 잔여 EMV 0.24 → 0.36. R-02 P 40 → 35(협의회 관계 개선, 잔여 0.60 → 0.53). 순증 +0.05 → 잔여 위협 EMV 합 7.26 → 7.31 (CCB 승인 시 등록부 v1.2에 반영) & ⑬ 항목 4 \\ \hline
계약/편익 & 변경 대가 \textbf{증액 0.19억}(변경협의체 09-30 합의). \textbf{B-08(38h → 12h)}: 대안 (a)(c)에서 유지, (b)는 14h로 목표 재정의(오정근·재무본부 판정). \textbf{타 프로젝트·운영 영향: 물류본부 야간조 편성(한동석), 3PL 배송 출발 시각(대성) → 프로그램 운영위 보고}. 협의회 회신 필요 & 헌장 §3 성공 기준 “가맹 발주 리드타임 12h, M+12 판정” \\ \hline}%
{범위 & 건수 0(변경) / FP +35, 누적 +245 & ±35건 / ±372 & 내 \\ \hline
일정 & A10 +1영업일 (임계, 이미 −10) & 컷오버 창 0 & \textbf{내(단독) — 그러나 ⑫ 회복 계획 전제} \\ \hline
원가 & +0.19 & 10.2 & 내 \\ \hline
품질 & TC +8 & — & 내 \\ \hline
리스크 & 잔여 EMV 합 순증 +0.05 & 8억 & 내 \\ \hline}%
{헌장 문언 변경: 아니오. 타 프로젝트·운영 영향: \textbf{예(물류본부·3PL)} → 규칙 ④ → \textbf{변경협의체(09-30 합의) → CCB(10-15) + 프로그램 운영위 보고(10-08)}. PM 전결 불가.}

% ---------------- 완성 예시: CR-10
\clearpage
\crsheet{CR-10 원가 허용범위 정의 변경 “+6\%(10.08억)” → “컨틴전시 전액 10.2억” (거버넌스 변경)}%
{CR 번호 / 제목 & CR-10 / 헌장 §12 PM 허용범위 표 원가 칸 및 거버넌스 §4 격자 원가·프로젝트 칸의 정의 확정 & 접수일 & 2026-09-18 (헌장 v1.1 수정 요청 4번 — G1에서 “프로그램 관리 규정 별표 해석 확인 후 확정”으로 유보 [추가 설정]) \\ \hline
요청자 / 접수 경로 & P1 PM (발견: Day2 원가 기준선 항목 9) / 결정자 정미란, 규정 해석 프로그램 PMO장(회신 2026-09-11) & 접수자 & P1 PMO \\ \hline}%
{\textbf{현재 문장:} 헌장 v1.1 §12 “원가 +6\% (10.08억)”[유보 표시], 회사 정본 §5.6 “원가 +6\% (P1 = 10.08억)”, 격자 확정본(Day2 §5.5 항목 7) “+10.2억 = 컨틴전시 전액 = 계획 원가의 +7.0\%” — 세 문서가 두 값을 갖는다.\par
\textbf{요청 문장:} 헌장 §12 원가 칸을 “\textbf{+10.2억 = P1 컨틴전시 전액 (계획 원가 145.8의 +7.0\%; 계획 원가 + 허용범위 = 집행 한도 156.0). 초과 예상 판정 기준 = EAC(⑫ 성과보고)가 집행 한도 156.0을 초과할 전망}”으로 확정. 거버넌스 §4 동일. 회사 정본 §5.6은 “프로그램 규정 초안 값(참고)”으로 각주.\par
\textbf{사유:} 프로그램 관리 규정 별표의 “+6\%”는 승인 예산(168.0)을 분모로 한 프로그램 계층 규정이고, P1의 PM 허용범위는 “PM이 스폰서 보고 없이 쓸 수 있는 돈”이어야 하며 그것이 컨틴전시 10.2다. 10.08로 두면 계획 원가 + 허용범위 = 155.88 $\ne$ 집행 한도 156.0이 되어 “허용범위 내인데 집행 한도 초과”가 논리적으로 가능하다. 프로그램 PMO장 회신(09-11): “별표의 6\%는 하한 권고이며 프로젝트별 컨틴전시 배정액으로 대체 가능.”\par
\textbf{안 바꾸면:} ⑫의 예외 보고 트리거가 두 개(155.88 / 156.0)가 되고, 윤태호 내부감사팀장의 “컨틴전시 집행 승인 근거” 질의에 두 문서가 다른 답을 한다.\par
\textbf{대안:} (a) \textbf{요청대로(컨틴전시 전액)} \quad (b) 10.08 유지, 컨틴전시 10.2 중 0.12는 스폰서 결재 \quad (c) 프로그램 규정 별표 자체를 개정 — 프로그램 계층 소관, P1 시한 밖.}%
{거버넌스·절차 변경 (FP 0) & 다음 CCB 10-15 — ⑫ 첫 EAC 기반 예외 판정(10-05 성과보고)이 이 정의를 쓴다 & IS-08(감리 지적 8번: 컨틴전시 집행 보고 누락)과 함께 대응 \\ \hline}%
{범위 & WBS·RTM·FP 영향 없음 & — \\ \hline
일정 & 없음 & — \\ \hline
원가 & 금액 이동 없음. 허용범위 값 10.08 → 10.20(+0.12). 예외 보고 트리거 “EAC > 156.0”으로 단일화. 컨틴전시 집행 절차(RACI 19행) 불변 & Day2 §4.5 항목 2·9 \\ \hline
품질 & 없음 & — \\ \hline
리스크 & 없음. 단, 리스크 허용범위(잔여 EMV 8억)와 원가 허용범위(10.2)의 관계를 헌장 각주로 명시(“EMV는 기대값, 컨틴전시는 자금 — 둘은 같지 않다”) & Day2 §5.6 해설 \\ \hline
계약/편익 & 계약 영향 없음(수행사 무관). 2020 내부감사 지적 “컨틴전시 집행 승인 근거”에 대한 답이 한 문장이 됨 & S-12 윤태호 \\ \hline}%
{\multicolumn{4}{|l|}{허용범위 자체의 정의 변경이므로 축별 판정은 해당 없음} \\ \hline}%
{헌장 §12 문언 변경: \textbf{예} → \textbf{CCB 승인(RACI 17행) + 스폰서 확인}, 프로그램 운영위 보고(규정 해석 첨부). 계약 서면: 불요.}

% ---------------- 완성 예시: CR-11
\clearpage
\crsheet{CR-11 데이터 전환 대상을 5년 이내 이력으로 한정, 5년 초과분 아카이브 (범위 축소)}%
{CR 번호 / 제목 & CR-11 / 헌장 §5 산출물 ⑤ 데이터 전환 대상 “1,842 테이블·7.4TB” → “5년 이내 이력(약 4.8TB), 5년 초과분(약 2.6TB) 아카이브 저장소 보관·미이관” & 접수일 & 2026-09-22 (헌장 v1.1 수정 요청 2번 — G1에서 “보존 요건 확인 후 확정” 조건부 유보) \\ \hline
요청자 / 접수 경로 & 박세연 IT본부장 · 재경팀장 (보존 요건 확인서 2026-09-22) / 원가 기준선 BL-01 조정표(데이터 전환 −2.0)의 근거 & 접수자 & P1 PMO \\ \hline}%
{\textbf{현재 기준선 문장:} 헌장 v1.1 §5 ⑤ “1,842 테이블·7.4TB 이관”[조건부 유보 표시]. 범위 기술서 v1.0 §2 ⑤ “5년 이내 이력 이관, 5년 초과분 아카이브(신규 가정)”. 원가 기준선 BL-01 조정표 “데이터 전환 11.40 → 9.40(−2.0, 범위 결정)”. 데이터 전환 용역 계약서는 \textbf{아직 11.4·7.4TB 기준}. SI FP 기준선에 5년 초과 이력 조회 화면·아카이브 조회 API 60 FP 포함.\par
\textbf{요청 문장:} 헌장 §5 ⑤를 범위 기술서 문장으로 확정. 용역 계약 변경(11.4 → 9.4, 대상 약 4.8TB·1,842 테이블 중 이력성 테이블 610의 5년 필터), SI 범위에서 5년 초과 이력 조회 화면 2종·API 1본 제외(FP −60), 아카이브 저장소(클라우드 저비용 계층, 연 0.06억 [추가 설정] P5 운영비) 접근 절차를 P5 운영 문서에 포함.\par
\textbf{사유:} 재경팀장 확인서 — 국세기본법상 장부·증빙 보존 5년, 상법상 상업장부 10년은 \textbf{SAP ECC 원장이 보존 주체}이며 OD-POS 거래 이력은 보조 기록으로 5년 초과분 원본 보관(아카이브)으로 충분. 세무조사 대응은 아카이브 원본 조회로 가능(재경팀 절차서 첨부). 5년 초과분 이관 시 전환 리허설 회당 +6시간(34시간 창 압박)과 검증 표본 대사 +40\% [추가 설정].\par
\textbf{안 바꾸면:} 원가 기준선 BL-01(9.4)과 용역 계약(11.4)의 2.0 불일치가 G2까지 남고, 리허설 1회차(2026-11-23)가 7.4TB로 실행되어 34시간 창 검증이 왜곡됨.\par
\textbf{대안:} (a) \textbf{요청대로} \quad (b) 7년 한정(세무 유예 감안) — 아카이브 1.2TB, 감액 −1.0, FP −30 \quad (c) 원안 유지(7.4TB 전량) — 관리예비비 2.0 복귀 요청, 리허설 시간 재산정.}%
{범위 축소 (헌장 §5 문언 확정, 계약 감액 2건) & 다음 CCB 10-15 — 매핑 종료 10-23(A09)·리허설 1회차 11-23 전 용역 계약 변경 서면 필요 & R-10(정본 판정 후 변경 — 5년 필터 규칙은 G2 정본 판정서에 포함), R-06(OD-POS 이력 구조 지식) \\ \hline}%
{범위 & WBS 1.5.7(매핑 대상 610 이력 테이블 5년 필터 규칙 추가), 1.5.9(리허설 데이터량), 1.2.2(주문 이력 조회 화면 2종 제외), 1.5.2(아카이브 조회 API 제외). RTM: SRS-905·906(HR-08 하위, 이력 조회) 삭제, SRS 신규 1(SRS-914 “5년 초과 이력은 아카이브 조회 절차로 제공”) → SRS 570 → 569. TC −9(TC-INT-036~040 삭제, TC-ORD-101~104 삭제), TC +2. \textbf{FP −60}(감액). 누적 FP 변동 +245 − 60 = \textbf{+185} & FP 산정서 부록 v1.1 §4; 범위 기술서 v1.0 §5 [신규] 가정 \\ \hline
일정 & A09(전환 매핑·정제, TF 5, 근접 임계) 150 → \textbf{140영업일}(5년 초과 테이블 정제 제외) → TF 5 → 15, 데이터 경로 여유 확대. A13(리허설 3회차) 회당 −6시간, 기간 불변. 임계경로 불변 & Day2 §3.5 CPM 표 A09·A11·A13; 용역사 PM 산정 2026-09-18 \\ \hline
원가 & SI \textbf{감액 −0.33억} = −60 × 0.0055. 용역 \textbf{−2.0}(BL-01 조정표 기반영 — 기준선 금액 변동 없음, 계약 금액만 11.4 → 9.4로 정합). 아카이브 저장소 연 0.06억은 P5 운영비(박준영 합의 09-24). 컨틴전시·관리예비비 이동 없음. \textbf{관리예비비 복귀 조건(조정표 “보존 요건이 이관을 요구하면 복귀”) 해제} & Day2 §4.5 항목 3 조정표 \\ \hline
품질 & TC −9 +2. 전환 검증 6종 중 “표본 대사” 모집단 7.4TB → 4.8TB(표본 수 불변). RTM 커버리지 불변. 감리 G2 문서 검토 대상에 아카이브 정책서 추가 & RTM v1.1 §10 \\ \hline
리스크 & R-10 잔여 불변(5년 필터 규칙을 정본 판정서 ④ 동결 규칙에 포함해 판정 후 변경 방지). R-06 잔여 0.15 유지. 신규 리스크 없음. \textbf{감액이므로 컨틴전시 소진에 영향 없음} & ⑬ \\ \hline
계약/편익 & 변경 대가 \textbf{감액 −0.33억}(SI, 변경협의체 09-30 합의 — 수행사 이의 없음), \textbf{용역 계약 변경 −2.0}(용역사 합의 09-25, 강현도). 편익 영향 없음(B-ID 무관). 타 프로젝트: P5(아카이브 접근 절차 운영 문서 12종에 포함) & 계약 변경 서면 2건(SI 변경 합의서 6호, 용역 변경 계약 1호) \\ \hline}%
{범위 & 건수 −1 / FP −60, 누적 +185 (+1.5\%) & ±35건 / ±372 (±3\%) & 내 \\ \hline
일정 & A09 TF 5 → 15 (개선) & — & 내 \\ \hline
원가 & −0.33 (감액), 용역 −2.0 계약 정합 & — & 내 \\ \hline
품질 & TC −7 순 & — & 내 \\ \hline
리스크 & 변동 없음 & 8억 & 내 \\ \hline}%
{헌장 §5 문언 변경: \textbf{예} → \textbf{변경협의체(09-30) → CCB(10-15)}, 계약 서면 2건. 편익 축 영향 없음 확인(오정근 서명 불요 — 편익 무관).}
\fi

% ================================================================== 8 CCB 의결서 (가로)
\landscapeon
\secbar{8. CCB 의결서 — 의결 · 조건 · 반대의견}
\guide{회의 정보(일자·정기/임시·출석·정족수·의장·배석), 안건별 의결(승인 / 조건부 승인 / 이연 / 기각)은 \textbf{택한 대안}으로 적는다. 조건에는 무엇·누가·언제까지·미이행 시 결과. 반대의견은 실명·사유로 적고 없으면 “없음”이라고 적는다(비워 두지 않는다). PM 전결 건의 사후 보고와 컨틴전시 집행 현황도 안건이다.}
{\footnotesize
\begin{tabularx}{\linewidth}{|L{28mm}|Y|}\hline
\hd{항목} & \hd{내용} \\ \hline
\B{\rd{8mm}}{}회의 & \Bg{[일자·시각, 정기/임시, 장소]}{2026-10-15(목) 14:00~16:10, 정기(셋째 목요일), 본사 7층 회의실 — 제10차} \\ \hline
\B{\rd{8mm}}{}출석 / 정족수 / 의장 & \Bg{[n/7 (실명) / 5 / 위원장]}{7/7 (정미란 위원장, 박세연, 오정근, 나영선, 최영주, 프로그램 PMO장, 재경팀장) / 5 / 정미란} \\ \hline
\B{\rd{8mm}}{}배석 & \Bg{[안건 설명·기록·수행사·계약]}{P1 PM(안건 설명), P1 PMO 기준선 담당(기록), 이재현(CR-08·09·11 변경 대가 확인, 의결 시 퇴장), 강현도(계약 서면), 한동석(CR-09 서면 의견 제출·불참)} \\ \hline
\B{\rd{10mm}}{}사후 보고 안건 & \Bg{[PM 전결 CR-nn — 이의 여부 / 컨틴전시 집행 현황]}{CR-07(인터페이스 미등재 2본 추가, FP +40, 0.22억, P1 PM 전결 09-17) — 이의 없음, 승인 확인. 컨틴전시 집행 현황(⑬ 항목 7): 누적 1.41 — \textbf{8월 CCB 보고 누락분(R-03 0.50) 소급 보고}(감리 지적 AF-1-08 대응)} \\ \hline
\end{tabularx}}
\vspace{2mm}
{\footnotesize
\begin{tabularx}{\linewidth}{|L{34mm}|L{48mm}|Y|L{52mm}|}\hline
\hd{안건} & \hd{의결 (택한 대안)} & \hd{조건 (무엇 · 누가 · 언제까지 / 미이행 시)} & \hd{반대의견 (실명 · 사유)} \\ \hline
\ifwbblank
\rd{18mm}CR-\g{[nn]} & \g{[승인 / 조건부 승인 / 이연 / 기각 — 대안 (x)]} & \g{① [ ] ② [ ] / 미이행 시 [ ]} & \g{[실명: 사유 / 없음]} \\ \hline
\rd{18mm}CR-\g{[nn]} & & & \\ \hline
\rd{18mm}CR-\g{[nn]} & & & \\ \hline
\rd{18mm}CR-\g{[nn]} & & & \\ \hline
\else
CR-08 파일럿 임시 연동 FP +90 & \textbf{승인 — 대안 (b)} (A12 편입, 0.495억 미배정 풀) & ① A12 편입 후 TF 5 이하로 떨어지면 P1 PM이 2영업일 내 보고(이재현) ② R-19를 등록부 v1.2에 정식 등재(P1 PM, 10-23) ③ 임시 연동 2본은 컷오버 후 30일 내 폐기·회수 절차를 컷오버 런북(Day4 ⑯)에 포함(박세연·박준영) ④ 파일럿 매장 재선정(IS-04) 결과를 P3가 11-13까지 통보 / 미이행 시 파일럿 가맹 3점 제외 & 없음 \\ \hline
CR-09 22:00 마감 & \textbf{조건부 승인 — 대안 (c)} (파라미터 구현 FP +35, 0.19억; 초기 운용 20:00) & ① 22:00 운용 전환은 물류본부 야간조 편성안(한동석)과 3PL 배송 출발 시각 합의를 첨부해 CCB 재상정 — 기한 2027-01 CCB(P2 안정화 검토 후) ② 협의회 회신(오정근, 10-23): “포털 선행 릴리스에 22:00 설정 기능 포함, 운용 시각은 물류 편성과 함께 2027-01 확정” ③ B-08 목표 12h는 유지, 14h 대안은 채택하지 않음 ④ R-16 잔여 EMV 갱신(P1 PM) / ①이 2027-01 CCB에 미상정 시 20:00 운용 확정 & \textbf{한동석(서면 대독)}: “야간조 편성은 P2 안정화 이후에나 검토 가능. 파라미터 구현에는 이의 없음.” — 반대의견으로 기재 \\ \hline
CR-10 원가 허용범위 정의 & \textbf{승인 — 대안 (a)} & ① 헌장 v1.1 §12 문언 확정·서명(정미란, 10-23) ② 거버넌스 §4 격자 확정본과 회사 정본 §5.6 각주(P1 PMO) ③ 프로그램 운영위 보고(프로그램 PMO장, 11월 운영위) ④ 윤태호 내부감사팀장 사본 송부 & 없음 \\ \hline
CR-11 전환 5년 한정 & \textbf{승인 — 대안 (a)} & ① 용역 변경 계약 1호(−2.0) 서명 10-23(강현도·용역사) ② SI 변경 합의서 6호에 CR-07·08·09·11 통합(+40 +90 +35 −60 = +105 FP, +0.5775억) 서명 10-23 ③ 5년 필터 규칙을 G2 정본 판정서 ④ 동결 규칙에 포함(박세연) ④ 아카이브 접근 절차를 P5 운영 문서 12종에 포함(박준영, G4 전) ⑤ 재경팀장 보존 요건 확인서를 헌장 §9 가정 부속으로 첨부 & 없음 \\ \hline
\fi
\end{tabularx}}
\B{\small\g{의결 후 조치: 기준선 BL-[ ] 발행 [일자] / 의결서 보관 위치 [ ] / 사본 [ ]}\par}{\small 의결 후 조치: 기준선 \textbf{BL-02} 발행 2026-10-23(항목 9). CCB 의결서는 형상 저장소 governance/ccb/2026-10-15.pdf, 윤태호 사본.\par}

% ------------------------------------------------------------------ 9 기준선 갱신
\clearpage
\secbar{9. 기준선 갱신 기록 (BL-0n)}
\guide{갱신 전후의 기준선 번호, FP 합계·계획 원가·PMB의 전후 값, 갱신되는 문서와 버전, 일자·서명자. \textbf{컨틴전시에서 기준선으로의 이동은 집행 한도를 바꾸지 않는다}는 것을 명시한다. FP를 고쳤으면 PV 곡선도 고친다 — 아니면 EV는 새 범위, PV는 옛 범위를 가리킨다. 감액 변경을 기준선에서 빼지 않으면 EV가 100\%를 넘는다.}
{\footnotesize
\begin{tabularx}{\linewidth}{|L{34mm}|L{44mm}|L{52mm}|Y|L{50mm}|}\hline
\hd{항목} & \hd{BL-\B{[전]}{01} (\B{[일자]}{G1 2026-04-30})} & \hd{BL-\B{[후]}{02} (\B{[일자]}{2026-10-23})} & \hd{변동} & \hd{근거} \\ \hline
\ifwbblank
\rd{8mm}FP 기준선 & & & \g{[CR별 델타 합]} & \g{[변경 로그 누적, FP 확인서 서명 3인]} \\ \hline
\rd{8mm}SI 계획 원가 & & & \g{[FP Δ × 단가]} & \\ \hline
\rd{8mm}계획 원가 (PMB) & & & & \g{[컨틴전시 → 기준선 이동]} \\ \hline
\rd{8mm}컨틴전시 잔여 & & & \g{[배정 이동; 집행은 별도]} & \\ \hline
\rd{8mm}집행 한도 & & & \g{[불변]} & \\ \hline
\rd{8mm}관리예비비 & & & & \\ \hline
\rd{8mm}PV 곡선 & & & \g{[갱신 월과 금액]} & \\ \hline
\rd{8mm}일정 기준선 & & & \g{[재기준선 여부]} & \\ \hline
\rd{8mm}범위 기술서 / WBS / 사전 & & & & \\ \hline
\rd{8mm}RTM & & & & \\ \hline
\rd{8mm}리스크 등록부 & & & & \\ \hline
\rd{8mm}헌장 & & & & \\ \hline
\else
FP 기준선 & 12,400 & \textbf{12,585} & +185 (CR-03 +22, CR-06 +58, CR-07 +40, CR-08 +90, CR-09 +35, CR-11 −60) & 변경 로그 누적, FP 확인서 BL-02(P1 PM·이재현·형상 담당 3인 서명) \\ \hline
SI 계획 원가 & 68.20 & 69.22 & +1.02 (185 × 0.0055 = 1.0175 → 1.02) & 변경 대가 누적 \\ \hline
계획 원가 (PMB) & 145.80 & \textbf{146.82} & +1.02 & 컨틴전시 → 기준선 이동 (미배정 풀) \\ \hline
컨틴전시 잔여 & 10.20 → 집행 후 8.79 & \textbf{7.77} & −1.02 (배정 이동; 집행 1.41은 별도) & ⑬ 항목 7 \\ \hline
집행 한도 & 156.00 & 156.00 & 불변 & 컨틴전시 내 이동 \\ \hline
관리예비비 & 12.00 & 12.00 & 불변 (CR-11로 복귀 조건 1건 해제) & — \\ \hline
PV 곡선 & 17개월, 2026-09 누적 77.83 & 2026-12 +0.33(CR-08·09 중 A12·S23 배분), 2027-01 +0.69, 누적 146.82 & 갱신 월만 & ⑫ 항목 3 각주 \\ \hline
일정 기준선 & v1.0 & v1.0 유지 (A12 25일·TF 5, A09 140일·TF 15 — 활동 내 흡수, 마일스톤·임계경로 불변) & 재기준선 아님 & ⑫ 회복 계획은 별도 \\ \hline
범위 기술서 / WBS / WBS 사전 & v1.0 & v1.1 (1.1.7·1.4.2·1.5.4·1.5.7·1.2.2 카드 갱신) & — & — \\ \hline
RTM & v1.1 (FZ-09) & v1.2 (StRS 215, SyRS 399, SRS 575, CR 이력 6건) & — & — \\ \hline
리스크 등록부 & v1.1 (G1) & v1.2 (⑬) & R-19 등재 & — \\ \hline
헌장 & v1.1 (§5·§12 유보 표시) & v1.1 확정본 (유보 해제, 서명 10-23) & — & — \\ \hline
\fi
\end{tabularx}}
\B{\small\g{서명: P1 PM \hrulefill\ 스폰서 \hrulefill\ 수행사 PM \hrulefill\ 형상 담당 \hrulefill}\par}{\small 서명(10-23): P1 PM / 정미란 / 이재현(FP 확인) / P1 PMO 형상 담당. 집행 한도 156.0 불변 — 스폰서가 “기준선이 늘었다 = 돈이 더 든다”로 읽지 않도록 항목 9에 명시.\par}

% ------------------------------------------------------------------ 10 변경 로그
\clearpage
\secbar{10. 변경 로그 (착수 이후 전 건)}
\guide{착수 이후 모든 CR의 한 줄 요약 — PM 전결 건과 기각·이연 건도 남긴다(CCB 안건만으로 만들면 누적이 틀린다). 누적 행은 분기 재확인(FP ±372)과 대조하는 용도이며 RTM 변경 이력 열과 번호가 같아야 한다.}
{\scriptsize\setlength{\tabcolsep}{2pt}
\begin{longtable}{|L{11mm}|L{15mm}|L{52mm}|L{16mm}|L{24mm}|C{10mm}|C{14mm}|L{24mm}|L{16mm}|L{28mm}|L{20mm}|L{12mm}|}\hline
\hd{CR} & \hd{접수일} & \hd{제목} & \hd{유형} & \hd{요청자} & \hd{FP Δ} & \hd{금액(억)} & \hd{자금원} & \hd{판정} & \hd{경로} & \hd{상태} & \hd{BL} \\ \hline \endhead
\ifwbblank
\rd{6mm}CR-01 & & & & & & & & & & & \\ \hline
\rd{6mm}CR-02 & & & & & & & & & & & \\ \hline
\rd{6mm}CR-03 & & & & & & & & & & & \\ \hline
\rd{6mm}CR-04 & & & & & & & & & & & \\ \hline
\rd{6mm}CR-05 & & & & & & & & & & & \\ \hline
\rd{6mm}CR-06 & & & & & & & & & & & \\ \hline
\rd{6mm}CR-07 & & & & & & & & & & & \\ \hline
\rd{6mm}CR-08 & & & & & & & & & & & \\ \hline
\rd{6mm}CR-09 & & & & & & & & & & & \\ \hline
\rd{6mm}CR-10 & & & & & & & & & & & \\ \hline
\rd{6mm}CR-11 & & & & & & & & & & & \\ \hline
\rd{6mm}CR-12 & & & & & & & & & & & \\ \hline
\rd{6mm}\textbf{누적} & & \g{[n건: 승인·조건부·기각]} & & & \g{[±n]} & \g{[ ]} & \g{[미배정 풀 소진]} & \g{[±372 대비]} & & & \\ \hline
\else
CR-01 & 2026-02-11 & 배달3사 인터페이스 3본 규격 v2 반영(벤더 규격 변경) & 오류 정정 & 박세연 & 0 & 0 & — & 내 & PM 전결 & 완료 & BL-01 \\ \hline
CR-02 & 2026-02-26 & 재고 정확도 리포트 두 분모 병기 추가(StRS-034) & 요구 변경 & 나영선·한동석 & +15 & 0.08 & 기준선 내(설계 스프린트 FP 배분 조정) & 내 & PM 전결 & 완료 & BL-01 \\ \hline
CR-03 & 2026-03-11 & 엑셀 발주 양식 3종 → 5종(StRS-072) & 요구 변경 & 오정근 & +22 & 0.12 & 미배정 풀 & 내 & PM 전결·CCB 사후 & 완료 & BL-02 \\ \hline
CR-04 & 2026-03-25 & 옴니 고도화 3종 R5 이연 확정(StRS-053 등) & 릴리스 이연 & 오정근 & 0 & 0 & — & 헌장 §13 & CCB(04-16) & 승인 & BL-01 \\ \hline
CR-05 & 2026-03-25 & 배달 주문 가맹 배분·정산(StRS-054) & 요구(기각) & 서정필 & 0 & 0 & — & Won't & Senior User & \textbf{기각} (기록 유지) & — \\ \hline
CR-06 & 2026-06-10 & FZ-05 접수 38건 처리 — 16건 승인, 22건 기각·헌장 수정 요청 & 요구 변경 & 영업본부(오정근 A) & +58 & 0.32 & 미배정 풀(재산정 공수 0.15는 R-07 배정) & 내 (R-07 발동) & 변경협의체 → PM 전결·CCB 사후 & 완료 & BL-02 \\ \hline
CR-07 & 2026-09-17 & 인터페이스 목록 갱신 — 폐기 2본 제외·미등재 2본 추가(IS-03) & 범위 정정 & 박세연 & +40 & 0.22 & 미배정 풀 & 내 & PM 전결·CCB 사후(10-15) & 승인 & BL-02 \\ \hline
CR-08 & 2026-09-16 & 파일럿 임시 재고 동기화 연동 & 범위 추가 & 정미란(수정 요청 1번) & +90 & 0.495 & 미배정 풀 & 내·헌장 §5 & 변경협의체 → \textbf{CCB} & 승인(조건 4) & BL-02 \\ \hline
CR-09 & 2026-09-17 & 가맹 발주 마감 20:00 → 22:00 & 요구 변경 & 서정필(S-10) & +35 & 0.19 & 미배정 풀 & 내·운영 영향 & 변경협의체 → \textbf{CCB} + 운영위 & 조건부 승인(파라미터) & BL-02 \\ \hline
CR-10 & 2026-09-18 & 원가 허용범위 정의 확정 & 거버넌스 & P1 PM(수정 요청 4번) & 0 & 0 & — & — & \textbf{CCB} + 스폰서 & 승인 & 헌장 v1.1 \\ \hline
CR-11 & 2026-09-22 & 전환 5년 한정·아카이브 & 범위 축소 & 박세연·재경팀장(수정 요청 2번) & −60 & −0.33 (+용역 −2.0) & 감액 & 내·헌장 §5 & 변경협의체 → \textbf{CCB} & 승인 & BL-02 \\ \hline
\textbf{누적} & & 11건: 승인 8 · 조건부 1 · 거버넌스 1 · 기각 1 & & & \textbf{+185} (+1.5\%) & \textbf{+1.02} (SI) / 용역 −2.0(BL-01 기반영) & 미배정 풀 소진 1.245 & ±372 이내 & & & \\ \hline
\fi
\end{longtable}}
\B{}{\small 검산: FP 누적 +22 +58 +40 +90 +35 −60 = +185; 12,400 × 0.03 = 372 → +185는 허용범위의 49.7\%. SI 금액 185 × 0.0055 = 1.0175 → 1.02. 미배정 풀 소진 = CR-03 0.12 + CR-06 0.32 + CR-07 0.22 + CR-08 0.495 + CR-09 0.19 − CR-11 0.33 → BL-02 편입 순 1.02(감액 반영 전 1.345, CR-11 감액은 기준선 차감). 2026-09-30 성과보고(⑫)는 BL-01 기준이므로 CR-03의 0.12는 “PV 없는 AC”로 별도 표시.\par}
\landscapeoff
\end{document}
